% !PRINTED 515
\weberpage{515}
%<<CONT
befriedigt, und zwar die erste oder die zweite, je nachdem
\weq{11}{\left(\frac{A}{m''}\right) = +1 \quad \text{oder} \quad \left(\frac{A}{m''}\right) = -1}
ist.

2. \gesperrt{Hieraus folgt, daß das Vorzeichen} $\left(\dfrac{A}{m''}\right)=\pm 1$ \gesperrt{nicht von der besonderen Form} $(A,B,C)$\gesperrt{, sondern nur von der durch diese Form repräsentierten Klasse abhängt, also ein Charakter der Formenklasse ist.}
Nun ist die Invariante einer zweiseitigen Klasse reell, und die Invarianten entgegengesetzter Klassen sind \gesperrt{konjugiert imaginär.} Für zwei entgegengesetzte Formen $(A,B,C)$, $(A,-B,C)$ ist aber das Vorzeichen (11) das gleiche, und daraus folgt, daß $\Phi(\mu,u)$ entweder reelle oder konjugiert imaginäre Wurzeln, und folglich \gesperrt{reelle Koeffizienten} hat, und sich daher nicht ändert, wenn $\mu$ durch den konjugiert imaginären Wert $\mu'$ ersetzt wird. Nun sind die beiden Fälle $m''\equiv 1$, $m''\equiv 3 \ (\mathrm{mod}\ 4)$ zu unterscheiden, weil es davon abhängt, welches Glied in (9) imaginär ist. Das eine Mal kommt $i$ nur in der Verbindung $i\sqrt{m'}$, das andere Mal in $i\sqrt{m''}$ vor, und demnach ist:
\weq{12}{\begin{aligned}
\Phi(\mu,u) = \tfrac{1}{2}\bigl[\Phi(\mu,u)+\Phi(\mu',u)\bigr] &= \Psi(\sqrt{m''},u): m''\equiv 1 \\
&= \Psi(\sqrt{m'},u): m''\equiv 3
\end{aligned}
\ (\mathrm{mod}\ 4).}
Die Funktion $\Psi$, und folglich auch $\Phi$, hat reelle Koeffizienten und enthält nur die eine der beiden Quadratwurzeln $\sqrt{m''}, \sqrt{m'}$. Die $\sqrt{m''}$ ist nach unserer Voraussetzung immer \gesperrt{irrational,} $\sqrt{m'}$ ist nur dann rational, wenn $m'$ ein Quadrat ist. Da $H(u)$ rationale Koeffizienten hat, so muß es, wenn $m''\equiv 1$ durch $\Psi(\sqrt{m''},n)$ und $\Psi(-\sqrt{m''},n)$, und wenn $m''\equiv 3$, und $m'$ kein Quadrat ist, durch $\Psi(\sqrt{m'},n)$ und $\Psi(-\sqrt{m'},n)$ teilbar sein.
% SOURCE-ORIGINAL-SIC: The registered page unambiguously prints n as the second argument in all four Psi terms, although equation (12) and H(u) indicate that u is mathematically expected.
Nun ändert $M$ sein Vorzeichen durch die gleichzeitigen Vorzeichenänderungen:
\weq{13}{
\begin{aligned}
(i,-i), (\sqrt{m''},-\sqrt{m''}): & m''\equiv 1 \\
(i,-i), (\sqrt{m'},-\sqrt{m'}): & m''\equiv 3
\end{aligned}
\ (\mathrm{mod}\ 4),}
und folglich ist, wenn nicht $m''\equiv 3$, und zugleich $m'$ ein Quadrat ist, in beiden Fällen:
\weq{14}{H(u) = \Phi(M,u)\,\Phi(-M,u),}
und diese Zerlegung ist nur dann nicht möglich, wenn $m''\equiv 1 \ (\mathrm{mod}\ 4)$ und zugleich $m'$ ein Quadrat oder $m'=1$ ist.
% SOURCE-ORIGINAL-SIC: The registered page prints m'' congruent to 1 here, although the sentence introducing (14) and the preceding argument indicate the exceptional case m'' congruent to 3.
%CONT>>
